\documentclass{article}
% Competition/workshop paper, non-anonymous (single-blind) camera-ready: the official style
% requires a track option together with [final]; [final] alone leaves the track name undefined.
\usepackage[sglblindworkshop,final]{neurips_2026}
\workshoptitle{Smart Buildings Challenge (SB26)}
\usepackage[utf8]{inputenc}
\usepackage[T1]{fontenc}
\usepackage{url}
\usepackage{booktabs}
\usepackage{amsfonts}
\usepackage{amsmath}
\usepackage{nicefrac}
\usepackage{microtype}
\usepackage{xcolor}
\usepackage{graphicx}
\usepackage{algorithm}
\usepackage{algpseudocode}
\usepackage{multirow}
\usepackage{hyperref}  % after float/algorithm so caption anchors are not duplicated
\input{results_macros}
% Figure include that degrades to a visible placeholder when the file is not built yet.
\newcommand{\figfile}[2]{\IfFileExists{#2}{\includegraphics[width=#1]{#2}}{\fbox{\parbox{#1}{\centering\textbf{??} missing figure \texttt{\detokenize{#2}}}}}}

\title{Schedule-Aware Demand-Based Reset: A Structured Rule-Based Controller for Cross-Building HVAC Transfer in the Smart Buildings Challenge}

\author{%
  {[Author Name]}\\
  {[Affiliation]}\\
  \texttt{neelsirivara@outlook.com}
}

\begin{document}

\maketitle

\begin{abstract}
The NeurIPS 2026 Smart Buildings Challenge asks for one HVAC control policy that operates
procedurally generated commercial buildings it has never seen, under 30-day weather strips
from five climates, trading occupant comfort against energy cost and carbon. We present
SADR (Schedule-Aware Demand-based Reset), a pure-numpy rule-based controller that reads the
plant topology from the building metadata and is otherwise built from facts shared by every
generated building: comfort is charged only while zones are occupied, fan power grows as
static pressure to the three-halves, the chiller has a constant coefficient of performance,
and all heat reaches a zone through a reheat coil whose discharge is capped at 308.15\,K.
SADR gates its mode on the occupancy count and schedule flags (with a self-tuning optimum
start), regulates each zone with a PI loop on \emph{net heat} translated into damper and
valve commands through the supply-air and discharge physics, and resets static pressure,
loop differential pressure and hot-water temperature from zone demand. On the 30 public
scenarios the round-4 version (v7) scores \RFConfPubCand{} on the leaderboard metric (mean absolute
per-step reward, lower is better) against \MetricPidWater{} for the strongest organizer
baseline (\texttt{pid\_water}) under identical conditions, a paired-bootstrap ratio of
\RFPwPubRatio{} (95\% CI \RFPwPubRatioLo--\RFPwPubRatioHi; \RFPwPubWins{} wins); on \NHeldAll{}
scenarios from \NHeldAllBuild{} generated buildings and unseen strips the ratio is
\RFPwOldRatio{} (CI \RFPwOldRatioLo--\RFPwOldRatioHi). Loss decompositions drove three
revisions: coldest supply air and a deeper setback (\CandWinnerRelRed{} on the public set);
zone loops scaled with floor area and box capacity for large zones (\HeldAllVFiveRelRed{} on
the held-out buildings, CI \HeldAllVFiveRelRedLo{} to \HeldAllVFiveRelRedHi); and, after
tracing the generated buildings' larger chiller term to the per-handler energy normalization
rather than to control, a pre-registered screen of discharge headroom, pressure-aware
requests, night air and set-back ventilation whose composite lowers the loss by
\RFConfPoolRed{} (CI \RFConfPoolRedLo{} to \RFConfPoolRedHi) on \RFConfPoolN{} held-out
scenarios no decision touched, by \RFConfFreshShipRed{} on fresh buildings in the shipped
buildings' size range and by \RFConfPubRed{} on the public set. On Kaggle's hidden scenarios v7
scores \RFiveKagVSeven, \RFiveKagGapPub{} times its public value; submissions of fixed earlier versions
and of the organizers' baseline show that the hidden buildings behave like the public ones and point to
harsher weather: the public scenarios plus the kit's unused Minneapolis station reproduce the probes'
ratios to v7. A fifth round adds a deeper weekend setback and lifts a request cap that bound in deep cold:
v8 lowers the loss by \RFiveConfPubRelAbs{} on the public set and \RFiveConfPubMinnRelAbs{} with Minneapolis
added (Kaggle: \RFiveKagVEight). We release the controller and report where it still loses.
\end{abstract}

\section{Introduction}

Learned building controllers have been studied for two decades, yet few run in real
buildings; the recurring obstacles are data-hungry training, fragile robustness and poor
generalization across buildings~\citep{wang2020rl}. The Smart Buildings Challenge (SB26)
isolates the last of these: a submission is a single policy class, instantiated on held-out
buildings drawn from the same procedural generator as six public ones and scored on 30-day
episodes under weather strips it has not seen~\citep{sb26kit}. The observation and action
vectors change shape with the building ($K$ air handlers and $V$ zones expose $K{+}1$ on/off
commands and $3K{+}2{+}2V$ continuous setpoints), and the simulator is Google's open-source
Smart Buildings simulator~\citep{goldfeder2023lightweight,goldfeder2025suite}.

We take the position that the right first entry in such a benchmark is a
\emph{structure-aware rule-based controller}: one that reads the plant topology from the
metadata and otherwise encodes physics every generated building shares. It transfers by
construction and shows which part of the score is reachable without learning.

\paragraph{Contributions.}
(i) An analysis of the SB26 reward and plant that yields the design facts above
(Section~\ref{sec:env}). (ii) SADR, a deterministic numpy controller (\PolicyMsOurs{}\,ms per
step) built from them (Section~\ref{sec:method}). (iii) A protocol of screens, ablation and
full-length loss decompositions that motivated three revisions, the last pre-registered and
confirmed on \RFFreshNBuild{} freshly generated buildings (Section~\ref{sec:results}). (iv) An
account of where the loss sits, why the controller lost on large zones, why the generated
buildings cost more, and where it still loses (Section~\ref{sec:analysis}), every number
generated from the result files.

\section{Environment and metric}
\label{sec:env}

\paragraph{Plant.}
Air handlers (AHUs) serve up to ten zones each through one VAV box per zone, mix return with
outside air (recirculation $r\le 0.7$) and can only cool: the supply is $\min(\text{mixed},
\text{SAT})$. Zone flow is $d\,F_{\max}\sqrt{P/P_{\text{design}}}$ for damper $d\in[0.1,1]$ and static
pressure $P\in[200,2000]$\,Pa; fan power is $\text{flow}\times P/\eta$, so it grows as $P^{3/2}$ once
dampers are open; the chiller's COP is a constant 3.5. Heat comes only from the VAV reheat coils on one
gas boiler loop (supply 310--350\,K, 0--20\,bar), whose effectiveness $\min(1, C_w/C_a,
(308.15-T_{\text{sa}})/(T_w - T_{\text{sa}}))$ caps the discharge at 308.15\,K; gas is billed as heat
delivered plus a tank loss.

\paragraph{Reward and metric.}
Each 5-minute step pays $0.6\,\text{(comfort shortfall)} + 0.2\,\text{(normalized energy
cost)} + 0.2\,\text{(normalized carbon)}$. The comfort shortfall is the occupancy-weighted
productivity loss (a deviation of $\Delta$\,K from the band costs $\Delta/5$ up to 5\,K and
$1+((\Delta-5)/5)^2$ beyond, normalized by total occupancy), so it is \emph{zero whenever
nobody is in the building}; occupants are present only on weekday comfort hours (arrivals
06--12, departures 13--19, none on US holidays) and only the building-wide count is
observed; the band is 294--297\,K in comfort hours and 289--298\,K otherwise. Energy is
normalized by a plant maximum that scales with $K$; in metric units one kW of electricity
costs about $2.1\times$ one kW of gas and one kW of chiller cooling about $0.6\times$. The
leaderboard metric is the mean over scenarios of
$|\text{mean per-step reward}|$, lower is better; it is dominated by the extreme climates
(the organizers' \texttt{pid\_water} loses \ExtremeRatioFairbanks$\times$ its mild-climate
mean in Fairbanks).

\paragraph{Baselines.}
The kit ships \texttt{bang\_bang}, \texttt{fixed\_setpoint}, \texttt{pid} (per-AHU PID on the
mean served-zone temperature, plant idle outside work hours) and \texttt{pid\_water}
(\texttt{pid} plus worst-zone tracking, a proportional economizer, proportional reheat and a
hot-water reset), the strongest and our main comparison; there is no learned baseline.

\section{Method: SADR}
\label{sec:method}

SADR is a single \texttt{TransferPolicy} class whose only state is per-zone integrals,
per-AHU dwell counters and the optimum-start statistics; all gains live in one parameter
dictionary (Table~\ref{tab:params}, Appendix~\ref{app:tuning}); Algorithm~\ref{alg:sadr} (Appendix~\ref{app:figs}) summarizes one step.

\paragraph{Topology from metadata.}
Index tables come from the observation field names, the positional AHU action lists, the
per-VAV action indices and the AHU-to-zone assignments; ranges and statistics from the
normalization metadata. Every action slot is written and clipped on every step; nothing depends
on the number of zones or handlers, the floor plan or the strip; non-finite inputs fall back to
schedule values; no files, no network, no randomness.

\paragraph{Modes.}
\textsc{Occupied} when the \texttt{comfort\_mode\_now} flag is set and the building is not
detected as closed; \textsc{Precondition} when unoccupied but comfort hours start within a
lead time (or \texttt{comfort\_mode\_soon}); \textsc{Setback} otherwise. ``Somebody is here''
is detected relative to the running minimum of the affinely normalized occupancy count; a
weekday on which nobody has arrived by \ParHolidayCheckH{}:00 is a holiday in the kit's
occupancy model and is treated as unoccupied until the count rises. The optimum-start lead is $\ell = \ell_0 +
\text{deficit}/(\hat{r}/s)$: the worst zone's distance from the comfort band over a
recovery-rate estimate $\hat{r}$ (K/h) re-estimated from each morning's recovery
(exponential average, weight \ParRateEma) with safety factor $s=\ParLeadSafety$, clipped to
$[\ParLeadMinH, \ParLeadMaxH]$\,h.

\paragraph{Zone loops on net heat.}
Each zone has heating and cooling targets, the sensor band shrunk by \ParMarginOcc\,K
(occupied) or \ParMarginPre\,K (preconditioning); two PI loops compute requested net heat
$q_h = k_q (T_h^\star - T)^+ + I_h$ and cooling $q_c = k_q (T - T_c^\star)^+ + I_c$ in watts
($k_q = \ParKqWPerK$\,W/K, integral gain \ParKiWPerKStep\,W/K per step); the integrals unwind
past the target, so a zone inside the band free-floats, and the larger request wins. Net
heat rather than flow because the zone-air node responds within minutes: a flow-proportional
controller with a 5-minute step oscillates, whereas a net-heat gain of the order of the air
node's conductance to the slab is stable and the integral carries the steady load.
\textsc{Setback} uses \emph{bursts} instead: a running handler pays the ventilation load of
every zone it serves, so the floor (eco band widened by \ParNightLowExt{}/\ParNightHighExt\,K)
is held by requesting \ParQSbW\,W from any zone that crosses it until it is \ParSbHystK\,K
inside, and the handler is off otherwise (holding the floor with the PI loops instead kept a
handler running for \VthreeAhuOnUnocc{} of the unoccupied steps on the coldest scenario).

\paragraph{From heat to actuators.}
A cooling zone needs capacity rate $F = q_c/(T - T_{\text{sa}})$, $T_{\text{sa}}$ being the
AHU's estimated supply temperature; a heating zone is bounded by the discharge limit, $F =
q_h/(308.15 - \delta - T)$. \emph{SAT:} a cooling handler commands the coldest value,
\ParSatCoolHi\,K: the coil load is $q_c + F\,(T_{\text{mixed}} - T)$ at a constant COP, so ventilation
term and fan power fall with flow (a Guideline-36 upward reset~\citep{ashrae2021g36,taylor2018vav} is
right in a real building, wrong here). A handler
that also heats is floored at \ParSatMixFloor\,K so its cooling zones are not overcooled into
reheat; a heating-only handler commands 305\,K so the coil idles. \emph{Static pressure} is
trim-and-respond in closed form, $P = P_{\text{design}}\max_z f_z^2$ clipped to
$[200, 2000]$\,Pa with $f_z$ each zone's requested fraction of design flow, so the most
demanding damper is fully open; dampers are $f_z/\sqrt{P/P_{\text{design}}}$. \emph{Reheat:} the discharge temperature that delivers $q_h$
at the delivered flow gives the water capacity rate $W_z = F_z\,(T_{\text{disch}} -
T_{\text{sa}})/(T_w - T_{\text{sa}})$; the loop differential pressure is reset so the largest
$W_z$ sits at \ParValveTarget{} of valve travel and each valve is $W_z/(k_v \sqrt{\Delta p})$.
An AHU runs iff a served zone requests more than \texttt{q\_on} (dwell \ParMinOnSteps{}
steps); the pump runs iff a valve is open. Recirculation is at the 0.7 cap except when a
cooling-only handler sees outside air colder than return air by more than the exhaust-fan
break-even ($0.00322$\,K/Pa$\times P$), when the outside-air fraction is solved so the mixed
air lands on the SAT target. The hot-water setpoint is 350\,K whenever heating was requested
in the last \ParHeatSeasonSteps{} steps and 310\,K otherwise: gas is billed as heat delivered,
so hotter water is strictly more useful at a tank dissipation of \TankLossExtraKw\,kW.

\paragraph{Scaling with the building (v6).}
Every VAV box in the generator has the same design flow (2\,m$^3$/s) and hot-water $k_v$, so
a 450\,m$^2$ zone has a fifth of the heating authority per square metre of a 100\,m$^2$ zone
and several times its thermal mass; v5's fixed 12\,kW cap, gains and building-wide lead were
tuned on the public zones (median \PubZoneAreaMedian\,m$^2$). v6 (i) reads each zone's floor
area from the metadata zone map (4-connected free cells in raster order reproduce the
simulator's room numbering) and scales $k_q$, $k_i$, the request cap and the setback burst with
area/\ParAreaRef\,m$^2$, clipped to $[\ParAreaScaleMin, \ParAreaScaleMax]$; (ii) bounds each
request by what the box delivers at a flow fraction $f_{\text{cap}}$ (\VSixFCapPre{} while
preconditioning or set back, \ParFCapOcc{} occupied) at the discharge limit, at 285\,K supply,
and on the water side (valve fully open at 20\,bar: asking for more raises the airflow, cools
the discharge and delivers \emph{less}); (iii) runs the optimum start per zone; (iv) raises the
setback floor of a zone whose learned rate could not recover it within the maximum lead. With
these off the code reproduces v5 bit for bit.

\paragraph{Round 4 (v7).}
\label{sec:v7}
The round-4 composite turns on five switches of the v6 code (all off reproduces v6 bit for bit),
each from a diagnosis in Section~\ref{sec:analysis}: (i) discharge margin $\delta$ from \VSixDischMargin{} to
\ParRFDischMargin\,K: the simulator caps the discharge itself, and at the cap the gas per joule of
net heat is $1 + 0.3\,(T - T_{\text{out}})/(308.15 - \delta - T)$, so headroom means less outside
air heated; (ii) an occupied zone may ask for flow above what its box delivers at the 200\,Pa
floor only in proportion to how far it is past its target (from \ParRFOccEzero\,K, full flow
\ParRFOccSpan\,K later): pressure above the floor costs fan power and enlarges every other zone's
forced minimum flow, while comfort is charged per occupant and a zone holds one or two whatever
its area; (iii) before a workday and outside heating season, set-back zones warmer than the
preconditioning target are flushed with outside air (floor pressure, coil off) while it is at
least \ParRFFlushDt\,K colder than the zone; (iv) in heating season the forced minimum flow of set-back
zones on a running handler is reheated to the zone temperature, at 1\,J of gas per joule instead of
the ventilated recovery's 1.9 later; (v) the preconditioning cap sits at the floor-pressure flow.

\paragraph{Tuning protocol.}
v4 came from a seeded random search (\NTuneConfigs{} configurations on six public scenarios at 1,440 steps:
defaults \TuneDefaultScreen, best \TuneBestScreen, \texttt{pid\_water} \TunePidWaterScreen; a one-at-a-time
ablation shows that doubling the integral gain to \ParKiWPerKStep\,W/K per step is the whole gain,
\AblRelDefAddKi), v5 from \NCandV{} physics-motivated candidates run at full length on the public set
(Appendix~\ref{app:tuning}). Rounds 3--5 screened single changes at 2,880 steps on public and
generated-building scenarios, combined the helpful ones and shipped only under a rule fixed beforehand (not
worse on the public set, better on held-out scenarios at full length); rounds 4 and 5 wrote the protocol down
before running anything (\texttt{PREREG.md}) and confirm on scenarios no decision touched: round 3 on
\NHeldAll{} scenarios of \NHeldAllBuild{} default-range buildings (\NHeldUnscr{} never screened), round 4 on
\RFFreshNBuild{} freshly generated buildings (three of its candidates were added after a first look at the
screen), round 5 on the shipped six under the one
weather station the kit ships but no public scenario uses (Minneapolis).

\section{Results}
\label{sec:results}

\paragraph{Setup.}
All runs use the kit's environment loader and \texttt{SimulationTracker}, 8,640 steps per scenario,
and the kit's \texttt{score\_leaderboard.py} formula with paired bootstrap CIs over scenarios. Rounds 1--4
ran on an Apple M5 Pro and a Windows laptop, round 5 on the laptop; the seeds are fixed and runs reproduce
across machines (SADR bit for bit, \texttt{pid\_water} within \MachineGapPidWater). Local \texttt{pid\_water} returns differ
from the shipped \texttt{baseline\_returns.json} by \GapMeanPidWater{} on average, so all comparisons use
local runs.

\begin{table}[t]
\caption{Leaderboard metric (mean over the \NScen{} public scenarios of $|$mean per-step
reward$|$; lower is better) and per-station means. Paired bootstrap 95\% CIs (\NBoot{}
resamples) are for SADR v6 minus the row (v7 against v6: Table~\ref{tab:r4}; v8 against v7: Section~\ref{sec:results}).
v8 ran on the laptop, the other rows on one machine (bit-identical across machines); the last row is the
organizers' shipped returns. Per-scenario numbers: Appendix~\ref{app:per_scenario}.}
\label{tab:main}
\centering\scriptsize
\setlength{\tabcolsep}{3pt}
\begin{tabular}{lcccccc c}
\toprule
policy & all & Fairbanks & Denver & Atlanta & Phoenix & Singapore & $\Delta$ vs v6 [95\% CI] \\
\midrule
SADR v8 (final, round 5) & \textbf{\RFiveConfPubCand} & \RFiveConfPubCandFairbanks & \RFiveConfPubCandDenver & \RFiveConfPubCandAtlanta & \RFiveConfPubCandPhoenix & \RFiveConfPubCandSingapore & Sec.~\ref{sec:results} \\
SADR v7 (round 4) & \RFConfPubCand &\RFConfPubCandFairbanks & \RFConfPubCandDenver & \RFConfPubCandAtlanta & \RFConfPubCandPhoenix & \RFConfPubCandSingapore & Table~\ref{tab:r4} \\
SADR v6 (round 3) & \MetricOurs & \MetricOursFairbanks & \MetricOursDenver & \MetricOursAtlanta & \MetricOursPhoenix & \MetricOursSingapore & -- \\
SADR v5 & \MetricOursVFive & \MetricOursVFiveFairbanks & \MetricOursVFiveDenver & \MetricOursVFiveAtlanta & \MetricOursVFivePhoenix & \MetricOursVFiveSingapore & \DiffOursVFive{} [\DiffLoOursVFive, \DiffHiOursVFive] \\
pid\_water & \MetricPidWater & \MetricPidWaterFairbanks & \MetricPidWaterDenver & \MetricPidWaterAtlanta & \MetricPidWaterPhoenix & \MetricPidWaterSingapore & \DiffPidWater{} [\DiffLoPidWater, \DiffHiPidWater] \\
fixed\_setpoint & \MetricFixedSp & \MetricFixedSpFairbanks & \MetricFixedSpDenver & \MetricFixedSpAtlanta & \MetricFixedSpPhoenix & \MetricFixedSpSingapore & \DiffFixedSp{} [\DiffLoFixedSp, \DiffHiFixedSp] \\
bang\_bang & \MetricBangBang & \MetricBangBangFairbanks & \MetricBangBangDenver & \MetricBangBangAtlanta & \MetricBangBangPhoenix & \MetricBangBangSingapore & \DiffBangBang{} [\DiffLoBangBang, \DiffHiBangBang] \\
pid\_water (organizer values) & \MetricPidWaterOff & \MetricPidWaterOffFairbanks & \MetricPidWaterOffDenver & \MetricPidWaterOffAtlanta & \MetricPidWaterOffPhoenix & \MetricPidWaterOffSingapore & \DiffPidWaterOff{} [\DiffLoPidWaterOff, \DiffHiPidWaterOff] \\
\bottomrule
\end{tabular}
\end{table}

\paragraph{Main result.}
SADR v8 (round 5, below) reaches \RFiveConfPubCand{} on the public set; v7 reaches \RFConfPubCand{} against
\MetricPidWater{} for \texttt{pid\_water}
(Table~\ref{tab:main}), a ratio of \RFPwPubRatio{} (95\% CI \RFPwPubRatioLo{} to
\RFPwPubRatioHi; wins \RFPwPubWins); v6 reached \MetricOurs{} (reduction \RelRedPidWater, CI
\RelRedLoPidWater{} to \RelRedHiPidWater; Wilcoxon $p\PWilcoxonPidWater$; per scenario in
Appendix~\ref{app:per_scenario}). On the public set v4 to v5 cut the metric by \CandWinnerRelRed{}
and v5 to v6 by \RelRedOursVFive{} (CI \RelRedLoOursVFive{} to \RelRedHiOursVFive), small because the
public zones are small.

\paragraph{Where the loss goes.}
Logging the plant energy rates every step splits the loss into comfort and the fan, chiller,
pump and boiler shares (Figure~\ref{fig:decomp}): on the public set v5 pays \DecPubVFiveGasShare{} for
gas, \DecPubVFiveCoolingShare{} for the chiller, \DecPubVFiveFanShare{} for fans and
\DecPubVFiveComfortShare{} for comfort (Fairbanks: gas \DecPubVFiveGasShareFairbanks), and
\texttt{pid\_water} \RegretSharePidWater{} for comfort.

\paragraph{Held-out buildings.}
Pooled over the \NHeldAll{} scenarios of the \NHeldAllBuild{} default-range generated buildings, run
once after the choice (Figure~\ref{fig:heldout}, Appendix~\ref{app:per_scenario}), v6 scores
\HeldAllMetricOurs{} against \HeldAllMetricOursVFive{} for v5 (difference \HeldAllVFiveDiff, CI
[\HeldAllVFiveDiffLo, \HeldAllVFiveDiffHi]; wins \HeldAllVFiveWins/\NHeldAll) and
\HeldAllMetricPidWater{} for \texttt{pid\_water} (ratio \HeldAllPidWaterRatio, CI
\HeldAllPidWaterRatioLo--\HeldAllPidWaterRatioHi); on the \NHeldUnscr{} never screened the reduction from
v5 is \HeldUnscrVFiveRelRed{} (CI \HeldUnscrVFiveRelRedLo{} to \HeldUnscrVFiveRelRedHi). The round-3 ledger
(Table~\ref{tab:candidates_r3}) keeps the failures (\NRThreeWorse{} of \NRThreeVariants{} variants worse,
among them v6 without the area scaling).

\paragraph{Round 4.}
Of the \RFNSingles{} single changes (Table~\ref{tab:candidates_r4}), \RFNHelps{} helped on the held-out
screen without hurting the public one, \RFNPubOnly{} on the public screen only, \RFNNs{} were not significant
and \RFNWorse{} were worse (caps below the floor-pressure flow; the night flush in heating season). At full
length (Table~\ref{tab:r4}) the composite v7 beats v6 on every set, by \RFConfPubRed{} on the public set (CI
\RFConfPubRedLo{} to \RFConfPubRedHi), \RFConfOldCleanRed{} on the clean old held-out scenarios and
\RFConfFreshRed{} on the fresh ones (\RFConfFreshShipRed{} in the shipped six's ranges), mostly gas
(Fairbanks \RFDecPubFairbanksRel), while comfort rises slightly: a few per cent, not a new regime.

\input{round5_kaggle_table}

\paragraph{Kaggle.}
On the Kaggle public leaderboard v7 scores \RFiveKagVSeven{} (rank \RFiveKagRankVSeven/\RFiveKagTeamsVSeven{}
when scored), \RFiveKagGapPub{} times its public-set number, although Kaggle's own run of a public scenario
matches the local total to 0.02\%: the gap is in the hidden scenarios (six buildings from the same generator
under five weather stations absent from the public set; about ten of them form the public split). Round 5
used Kaggle as an oracle: \RFiveNProbes{} submissions of fixed earlier files (v6, v5 and the organizers'
\texttt{pid\_water}) and a pre-registered comparison of every local set on the level of v7 and on the
shape $E$, how well a set reproduces each policy's ratio to v7 (Table~\ref{tab:kaggle}). Two probes
discriminate. v5's ratio to v7 on Kaggle (\RFiveKagRatioVFive) is the public set's (\RFivePubRatioVFive),
not that of the larger default-range buildings (\RFiveOldRatioVFive), whose large zones v6 fixed: the
hidden buildings behave like the shipped six. \texttt{pid\_water}'s ratio (\RFiveKagRatioPid) falls in cold
weather, where both controllers pay the same ventilation physics (public stations \RFivePubRatioPid,
Minneapolis \RFiveMinnRatioPid), and is reproduced by the public set with the Minneapolis scenarios added
(\RFivePubMinnRatioPid; $E=\RFivePubMinnShapeErr$ against \RFivePubShapeErr{} for the public set alone):
the gap points to weather, not buildings. Even on that set v7 reaches only \RFivePubMinnLevel{} of its Kaggle
score, so the hidden stations are harsher still. A public notebook shows the same gap for three controllers built on
\texttt{pid\_water} ($\times$\RFiveThirdARatio--\RFiveThirdCRatio), and its baseline's Kaggle score
(\RFiveThirdAKag) is close to our \texttt{pid\_water} submission (\RFiveKagPid). v6 scores \RFiveKagVSix{}
(\RFiveKagRatioVSix{} of v7, as on the local sets), a difference below the noise of a ten-scenario split.
Read literally, the rule names the fresh default-range set FD, whose $E$ rests on the v6 ratio alone, the one
that every set matches; round 5 therefore required its candidate to win on both FD and P+M.

\paragraph{Round 5.}
Round 5 screened \RFiveNSingles{} single changes to v7 at \RFiveScreenSteps{} steps on eleven held-out
scenarios (six of them the shipped buildings under Minneapolis strips outside set M, one the public
building~7 in the January Fairbanks strip) and four public ones (Table~\ref{tab:candidates_r5}):
\RFiveNHelps{} helped, \RFiveNNs{} were not significant and \RFiveNWorse{} were worse (a longer lead before
Mondays, \RFiveMonHeldRelAbs; recovering at the flow cap from the start, \RFivePreburstHeldRelAbs; faster
setback bursts); sensor-lag compensation did nothing measurable. Two families helped: a deeper
setback far from the next workday (floor 6\,K lower and ceiling 4\,K higher, ramped out 10--22\,h before the
start; \RFiveWkHeldRelAbs{} lower on the held-out screen, \RFiveWkPubRelAbs{} on the public one) and a request
cap of 20\,kW per 100\,m$^2$, which changed \RFiveQcapNChanged{} of the eleven held-out scenarios and lowered
building~7's January loss by \RFiveQcapBSevenRelAbs. Their composite beat v7 by \RFiveChosenHeldD{} (CI [\RFiveChosenHeldLo,
\RFiveChosenHeldHi], wins \RFiveChosenHeldWins) and \RFiveChosenPubD{} (CI [\RFiveChosenPubLo,
\RFiveChosenPubHi]) $\times10^{-5}$ per step on the two screens. At full length it lowers the loss by \RFiveConfPubRelAbs{} on the public set (CI
\RFiveConfPubRelHiAbs{} to \RFiveConfPubRelLoAbs; wins \RFiveConfPubWins), by \RFiveConfMinnRelAbs{} on M (CI
\RFiveConfMinnRelHiAbs{} to \RFiveConfMinnRelLoAbs; it changes \RFiveConfMinnNChanged{} winter scenarios and wins all of
them) and by \RFiveConfPubMinnRelAbs{} on P+M (CI \RFiveConfPubMinnRelHiAbs{} to \RFiveConfPubMinnRelLoAbs), but on FD
only by \RFiveConfFreshDefRelAbs{}: the CI's upper end is below zero for the pre-registered bootstrap seed and for
\RFiveSeedsFreshDefBelow{} of 10 seeds in all, so FD is a tie (a few hot-strip scenarios are up to
\RFiveConfFreshDefWorstRelAbs{} worse; the gain is in the cold strips). The rule is met as written and v8 is the
submission. On Kaggle it scores \RFiveKagVEight{} against v7's \RFiveKagVSeven, within the noise of the public split.

\section{Analysis}
\label{sec:analysis}

\paragraph{Why it wins, and where the loss sits.}
\texttt{pid\_water} keeps \OccOutPidWater{} of occupied zone-steps outside the band (one loop per handler,
a step start at 06:00; Figure~\ref{fig:hours}); SADR v6 keeps \OccOutOurs{} and uses far less energy
(floor pressure with wide dampers, least reheat flow at the discharge cap, no conditioning while nobody
is charged for comfort). For v6, \UnoccShareOurs{} of the loss is incurred on unoccupied steps, most at
hour \PeakLossHourOurs{}, the warm-up before 06:00. Fairbanks is gas: with a mandatory 30\% outside-air
fraction at $-20^\circ$C heat costs about $1.9$\,J of gas per joule at the discharge cap; Phoenix and
Singapore are chiller electricity, where the coldest supply at the lowest flow is cheapest.

\paragraph{What rounds 3 and 4 found.}
v5's 350--430\,m$^2$ zones sat 2--4\,K below the band until noon in the cold strips (fixed cap, single lead);
v6 cuts the largest-quartile recovery from \DecHeldRVFiveRecoveryBig\,h to \DecHeldROursRecoveryBig\,h. The
generated buildings' larger chiller term (\DecHeldROursCooling{} against \DecPubOursCooling{} per step) is
not a control defect: per square metre they cool no more (\RFCoolWmRThree{} against \RFCoolWmPub\,W/m$^2$),
but energy is normalized per handler and they have \RFAreaPerAhuRThree{} instead of \RFAreaPerAhuPub\,m$^2$
of zone per handler. In the cold strips outside-air heat is \RFVentGasPub{} of the gas. The kit's default
generator ranges do not reproduce the shipped six (seed 0: \RFSeedZeroDefaultSize{} ft against
\RFTrainOneSize).

\paragraph{What round 5 found.}
Under the Minneapolis strips (set M) v7 pays \RFiveDecMinnGasShare{} of its loss for gas and
\RFiveDecMinnComfortShare{} for comfort, \RFiveCapComfortShareOfMinn{} of it in one scenario
(\texttt{building\_train\_7} in January): its \RFiveCapNZones{} worst zones sit \RFiveCapDefMin--\RFiveCapDefMax\,K
below the band on \RFiveCapStepsMin--\RFiveCapStepsMax{} occupied steps with dampers open at
$\le$\RFiveCapPaMax\,Pa, because the request cap (\RFiveCapReqKwMin--\RFiveCapReqKwMax\,kW, a v5 constant
scaled by area) binds at about half of what box and coil deliver (\RFiveCapWaterKwMin--\RFiveCapWaterKwMax\,kW).
On the other sets the 06:00 hour carries \RFiveComfHSixPub--\RFiveComfHSixFresh{} of v7's comfort loss;
handlers start \RFiveAhuStartsMin--\RFiveAhuStartsMax{} times a day, so cycling is not a lever. Starter-kit
pull request~\#2 (unmerged) bills boiler gas from a tank energy balance and closes a gap that rewards
saw-toothing the hot-water setpoint (heat billed at 0.38 of its value); in a patched kit v7's loss falls by
\RFivePrTwoVSevenRelAbs{} on the \RFivePrTwoVSevenN{} public Fairbanks scenarios while \texttt{pid\_water}'s rises by
\RFivePrTwoPidRelAbs{}: SADR's score does not rest on that gap; the baseline's partly does.

\paragraph{Where it loses.}
SADR v7 is worse than \texttt{pid\_water} on \RFNScenWorseCand{} of \NScen{} public scenarios;
its largest residual loss is on \RFMaxCandScen{} (\RFMaxCandScore), almost all gas. What remains is the
morning recovery in the coldest strips, kept cheap rather than perfect (Figure~\ref{fig:hours}), and box
capacity on the largest zones; a learned policy's advantage should be looked for in the recovery.

\section{Limitations}
SADR encodes simulator constants (discharge limit, valve $k_v$, constant COP) that a real building would
commission, where cold SAT would be weighed against humidity and a weekend floor near freezing against frost
protection. Selection used the public set and screens, so the public number is optimistic; Kaggle's hidden
stations are harsher than any local set. MPC or a learned residual on top of SADR is the next step.

\section{Reproducibility and disclosure}
Code: \url{https://github.com/neelsirivara/neurips2026-smart-buildings}; \texttt{submission/transfer\_policy.py} is
the policy, \texttt{experiments/src/} generates every number (\texttt{finalize.sh}); seeds fixed; CPU only. An LLM
coding assistant (Claude) wrote code and drafted this text; every number is checked in \texttt{paper/AUDIT.md}.

\bibliographystyle{plainnat}
\bibliography{refs}

\clearpage
\appendix
\section{Per-scenario results}
\label{app:per_scenario}
Tables~\ref{tab:per_scenario}, \ref{tab:per_scenario_held} and~\ref{tab:per_scenario_held_r3}
list every scenario's total return and leaderboard score (SADR v6), and
Table~\ref{tab:per_scenario_fresh} the round-4 confirmation set, Table~\ref{tab:r4} round 4 at full length and
Table~\ref{tab:per_scenario_r5m} the round-5 Minneapolis set; all values are generated from the
result files by \texttt{make\_macros.py}.
\begin{table}[h]
\caption{Round 4 at full length (8,640 steps): the composite v7 (\RFChosen) against the shipped
v6 on the same scenarios; $\Delta$ = paired mean difference in $10^{-5}$ per step with its
bootstrap 95\% CI (\NBoot{} resamples; negative favours v7). Old clean = the round-2/3 held-out
scenarios no screen touched; fresh = \RFFreshNBuild{} buildings generated for round 4 and run only
for this table, half in the shipped six's parameter ranges, half in the kit's default ranges.}
\label{tab:r4}
\centering\scriptsize
\setlength{\tabcolsep}{4pt}
\begin{tabular}{lrrrrrr}
\toprule
set & $n$ & v6 & v7 & $\Delta$ [95\% CI] & relative [95\% CI] & wins \\
\midrule
public & \RFConfPubN & \RFConfPubVSix & \RFConfPubCand & \RFConfPubDiff{} [\RFConfPubLo, \RFConfPubHi] & \RFConfPubRel{} [\RFConfPubRelLo, \RFConfPubRelHi] & \RFConfPubWins \\
old held-out, clean & \RFConfOldCleanN & \RFConfOldCleanVSix & \RFConfOldCleanCand & \RFConfOldCleanDiff{} [\RFConfOldCleanLo, \RFConfOldCleanHi] & \RFConfOldCleanRel{} [\RFConfOldCleanRelLo, \RFConfOldCleanRelHi] & \RFConfOldCleanWins \\
fresh, shipped ranges & \RFConfFreshShipN & \RFConfFreshShipVSix & \RFConfFreshShipCand & \RFConfFreshShipDiff{} [\RFConfFreshShipLo, \RFConfFreshShipHi] & \RFConfFreshShipRel{} [\RFConfFreshShipRelLo, \RFConfFreshShipRelHi] & \RFConfFreshShipWins \\
fresh, default ranges & \RFConfFreshDefN & \RFConfFreshDefVSix & \RFConfFreshDefCand & \RFConfFreshDefDiff{} [\RFConfFreshDefLo, \RFConfFreshDefHi] & \RFConfFreshDefRel{} [\RFConfFreshDefRelLo, \RFConfFreshDefRelHi] & \RFConfFreshDefWins \\
fresh, all & \RFConfFreshN & \RFConfFreshVSix & \RFConfFreshCand & \RFConfFreshDiff{} [\RFConfFreshLo, \RFConfFreshHi] & \RFConfFreshRel{} [\RFConfFreshRelLo, \RFConfFreshRelHi] & \RFConfFreshWins \\
clean held-out, pooled & \RFConfPoolN & \RFConfPoolVSix & \RFConfPoolCand & \RFConfPoolDiff{} [\RFConfPoolLo, \RFConfPoolHi] & \RFConfPoolRel{} [\RFConfPoolRelLo, \RFConfPoolRelHi] & \RFConfPoolWins \\
\bottomrule
\end{tabular}
\end{table}
\input{per_scenario_table}
\input{per_scenario_table_r4}
\input{per_scenario_table_r5}
\begin{figure}[h]
\centering
\figfile{0.92\linewidth}{figures/fig_round4_buildings.pdf}
\caption{Round 4: change of the metric from v6 to v7 on every generated held-out building
(diamonds: the building's five strips pooled; dots: single scenarios; negative is better), each
run once at full length; buildings grouped by set and ordered by zone count.}
\label{fig:r4buildings}
\end{figure}
\begin{figure}[h]
\centering
\figfile{0.86\linewidth}{figures/fig_scores_by_scenario.pdf}
\caption{Per-scenario leaderboard score (log scale, lower is better), SADR v6 vs
\texttt{pid\_water} on the same machine, grouped by station and labelled by building (b =
public, g = generated held-out). Left: the 30 public scenarios; right: the \NHeld{}
held-out scenarios of buildings 0501 and 0502.}
\label{fig:scores}
\end{figure}

\section{Parameters, ablation and candidates}
\label{app:tuning}
\begin{table}[h]
\caption{Final SADR parameters (v8); values that changed from v4 to v5 are marked; the fifth row
is new in v6 (round 3), the discharge margin (v6: \VSixDischMargin\,K), the preconditioning cap
(v6: \VSixFCapPre) and the seventh row are round 4, the last row round 5. The full dictionary is in the code.}
\label{tab:params}
\centering\footnotesize
\begin{tabular}{llll}
\toprule
occupied margin & \ParMarginOcc\,K & precondition margin & \ParMarginPre\,K \\
setback widening (low/high) & \ParNightLowExt{}/\ParNightHighExt\,K (v4: \VFourNightLowExt/\VFourNightHighExt) & $k_q$, $k_i$ & \ParKqWPerK, \ParKiWPerKStep\,W/K \\
cooling SAT / mixed floor & \ParSatCoolHi{}/\ParSatMixFloor\,K (v4 ceiling: \VFourSatCoolHi) & discharge margin & \ParDischMarginK\,K \\
lead safety / rate$_0$ & \ParLeadSafety{} / \ParRateZeroKph\,K/h & lead range & \ParLeadMinH--\ParLeadMaxH\,h \\
heating-season memory & \ParHeatSeasonSteps{} steps & AHU dwell & \ParMinOnSteps{} steps \\
area scaling (v6) & per \ParAreaRef\,m$^2$, clip $[\ParAreaScaleMin, \ParAreaScaleMax]$ & flow cap $f_{\text{cap}}$ (v6) & \ParFCapPre{} pre/setback, \ParFCapOcc{} occ. \\
occupied pressure cap (v7) & from \ParRFOccEzero\,K, span \ParRFOccSpan\,K & night flush (v7) & $\ge$\ParRFFlushDt\,K, $\le$\ParRFFlushWindow\,h before work \\
weekend setback (v8) & $-$\RFiveParWkLo/$+$\RFiveParWkHi\,K, out \RFiveParWkRampEnd--\RFiveParWkRampStart\,h before work & request cap (v8) & \RFiveParQcapKw\,kW/100\,m$^2$ (v7: \RFiveParQcapKwVSeven) \\
\bottomrule
\end{tabular}
\end{table}

Table~\ref{tab:ablation} reverts each tuned v4 value on its own (and adds each on its own to
the defaults) on the screening set; Table~\ref{tab:candidates} lists every v5 candidate at
full length on the 30 public scenarios against v4 on the same machine;
Table~\ref{tab:candidates_r3} lists every round-3 (v6) variant on both screening sets against
v5, Table~\ref{tab:candidates_r4} every round-4 candidate against v6 and Table~\ref{tab:candidates_r5}
every round-5 candidate against v7. The ideas tried and rejected, with their numbers, are in
\texttt{experiments/results/round3/candidates.json}, \texttt{round4/candidates.json} and
\texttt{round5/screen\_summary.json}.
\input{tuning_tables}
\input{tuning_tables_r5}

\section{Algorithm and additional figures}
\label{app:figs}
\begin{figure}[h]
\centering
\begin{minipage}[t]{0.5\linewidth}\centering
\figfile{\linewidth}{figures/fig_loss_decomp.pdf}
\caption{Per-station decomposition of the per-step loss (metric units) into comfort, fan,
chiller, pump and gas on the 30 public scenarios, SADR v5 and v6.}
\label{fig:decomp}
\end{minipage}\hfill
\begin{minipage}[t]{0.44\linewidth}\centering
\figfile{\linewidth}{figures/fig_heldout_buildings.pdf}
\caption{Metric per generated held-out building (log scale, lower is better; ordered by zone
count): \texttt{pid\_water}, SADR v5 and v6, each run once at full length on the building's
five strips.}
\label{fig:heldout}
\end{minipage}
\end{figure}
\begin{algorithm}[h]
\caption{SADR, one step (vectorized over zones $z$ and handlers $k$)}
\label{alg:sadr}
\begin{algorithmic}[1]
\State mode $\gets$ \textsc{Occupied} / \textsc{Precondition} / \textsc{Setback} from flags, time, occupancy count, lead $\ell$
\State $(T_h^\star, T_c^\star)_z \gets$ band for the mode; update $I_h, I_c$; $q_h, q_c \gets$ PI requests (W), or setback bursts
\State run$_k \gets$ hysteresis on $\max_{z\in k}(q_h, q_c)$
\State SAT$_k \gets$ \ParSatCoolHi\,K if cooling (\ParSatMixFloor\,K if also heating), 305\,K otherwise; recirc$_k \gets$ economizer
\State $F_z \gets q_c/(T_z - T_{\text{sa},k})$ or $q_h/(308.15-\delta-T_z)$; $P_k \gets P_{\text{design}}\max_z f_z^2$; damper$_z \gets f_z/\sqrt{P_k/P_{\text{design}}}$
\State $W_z \gets F_z (T_{\text{disch},z} - T_{\text{sa}})/(T_w - T_{\text{sa}})$; $\Delta p \gets (\max_z W_z / (0.9 k_v))^2$; valve$_z \gets W_z/(k_v\sqrt{\Delta p})$
\State $T_w^{\text{sp}} \gets 350$\,K if heating season else 310\,K; pump $\gets$ any valve open
\State scatter all setpoints into the normalized action vectors; clip
\end{algorithmic}
\end{algorithm}
\begin{figure}[h]
\centering
\figfile{\linewidth}{figures/fig_loss_by_hour.pdf}
\caption{Mean per-step loss by hour of day over the \NAnalysis{} public scenarios, split
into comfort regret and energy\,+\,carbon; comfort hours (06--19) shaded. Left: SADR v6;
right: \texttt{pid\_water}.}
\label{fig:hours}
\end{figure}
\begin{figure}[h]
\centering
\figfile{\linewidth}{figures/fig_trace_building_train_8_fairbanks_8_day8.pdf}
\caption{One weekday on \texttt{building\_train\_8/fairbanks\_8}: zone-temperature
envelope (min--max shaded, median line) against the comfort band (green wash), with the
building occupancy count shaded in grey. Left: SADR v6; right: \texttt{pid\_water}.}
\label{fig:trace}
\end{figure}
\begin{figure}[h]
\centering
\figfile{0.5\linewidth}{figures/fig_tuning.pdf}
\caption{Screening score (six scenarios, five days) of every random-search configuration,
sorted, with the design defaults highlighted and \texttt{pid\_water} on the same runs.}
\label{fig:tuning}
\end{figure}

\end{document}
